\section{Conclusiones}\label{sec:conclusion}

He estudiado cómo el aumento del salario mínimo peruano de 2022, de 930 a 1025 soles, afectó a trabajadores con distintos niveles educativos, usando datos trimestrales de la ENAHO y un diseño de diferencias en diferencias que aprovecha la exposición marcadamente distinta de los trabajadores de baja y alta educación a una única reforma nacional. Mi hallazgo central es que la reforma desplazó la composición del empleo poco calificado y no el nivel de su remuneración: el empleo formal cayó 4.6 puntos porcentuales, una quinta parte de su nivel de base, y el trabajo independiente aumentó 3.6 puntos, con tendencias previas de la formalidad que superan su prueba conjunta y una deriva residual que cuantifico y corrijo, desplazamientos mayores donde el piso cortó más hondo en la distribución salarial regional y la confirmación de un diseño de exposición continua con inferencia por aleatorización. Los datos de panel enlazados ubican el margen: los trabajadores formales en su puesto fueron retenidos, mientras que la entrada a empleos formales de los trabajadores de baja educación cayó casi a la mitad, de modo que el piso operó cerrando la puerta a nuevos vínculos formales y no disolviendo los existentes. Ninguna respuesta salarial acompañó el desplazamiento: los salarios relativos de los trabajadores de baja educación no subieron ni en la media ni en ningún punto de la distribución, y la proporción de asalariados que gana por debajo del piso no cayó, lo que señala el incumplimiento como la razón por la que el piso no pudo elevar la remuneración. El efecto sobre el empleo total no puede precisarse y se confunde con la recuperación pospandemia desigual. En torno a la reforma de enero de 2025, de igual magnitud, ni los salarios ni la composición respondieron. Examinar directamente el contraste favorece una lectura basada en la posición del piso frente a una basada en la rotación del mercado laboral: el aumento de 2022 barrió una franja que concentraba la cuarta parte de la remuneración formal de baja educación, mientras que el de 2025 barrió una franja más delgada en una distribución que el piso ya había rebasado, aun cuando la creación de vínculos formales era, si acaso, mayor en 2024. El margen de reasignación opera cuando el nuevo piso atraviesa el rango donde efectivamente se ubica la remuneración formal, y queda inactivo cuando ya lo ha superado.

La lectura más económica de estos resultados es que, en un mercado laboral donde la informalidad es generalizada y la fiscalización débil, un aumento moderado del mínimo legal, erosionado por la inflación, no eleva la remuneración de los poco calificados; lo que mueve, cuando mueve algo, es la asignación de los trabajadores entre arreglos formales e informales. Esto matiza la visión del efecto faro del salario mínimo peruano que surge de los estudios de reformas anteriores y apunta al margen informal como el canal central de ajuste.

De aquí se desprenden dos líneas de investigación futura. La primera es afinar la evidencia sobre transiciones con datos enlazados de mayor frecuencia, ya sea la EPEN mensual o los registros administrativos de la Planilla Electrónica, que permitirían fechar dentro del año la respuesta por el margen de las contrataciones y seguir a las mismas empresas y trabajadores a lo largo de ambas reformas; mi panel anual establece el margen, pero no su dinámica fina. La segunda es estudiar directamente la interacción entre la política de salario mínimo y la fiscalización, pues mis resultados sugieren que los réditos de elevar el piso dependen de que este pueda hacerse vinculante. A medida que se acumulen más trimestres posteriores a la reforma, también será posible rastrear el ajuste de largo plazo que una ventana corta como la mía no logra capturar.
